\documentclass[12pt]{article}
\usepackage[utf8]{inputenc}
\usepackage[T1]{fontenc}
\usepackage{hyperref}
\title{Yetfa Language Resources for K-12 Curriculum Development}
\author{Compiled by Claude Sonnet 5.5 (claude-sonnet-5-5) for Andr\'as Kornai}
\date{}
\begin{document}
\maketitle

\section{Introduction}
Yetfa (Biksi-Yetfa) is a Papuan language of Jetfa District, Pegunungan Bintang Regency, Highland Papua, Indonesia. It extends into Papua New Guinea. Wikipedia cites about 1,000 speakers (1996). Another source gives about 1,400. Its classification is unsettled: Foley (2018) treats it as an isolate, and earlier work placed it in the Pauwasi group.

\section{Primary Educational Resources}
The searches found no primers, textbooks or school programmes for Yetfa.

\section{Linguistic and Reference Materials}
\begin{itemize}
\item Wikipedia, Biksi-Yetfa language: \url{https://en.wikipedia.org/wiki/Biksi-Yetfa_language}
\item Numeral data for Yetfa (Max Planck Institute for Evolutionary Anthropology): \url{https://lingweb.eva.mpg.de/channumerals/Yetfa.htm}
\end{itemize}

\section{Religious and Community Resources}
None were found.

\section{Digital and Audio Resources}
The Endangered Languages Project has a Yetfa record: \url{https://endangeredlanguages.com/lang/3494}

\section{Curriculum Development Context}
Very little exists for this language. Any curriculum would need new documentation work.
\end{document}